% FONTE TEMA https://github.com/matze/mtheme
\documentclass[aspectratio=1610]{beamer}
%\documentclass[aspectratio=1610, handout]{beamer}
\usepackage[utf8]{inputenc}
\usepackage{ragged2e}
\usepackage{xcolor}
\usepackage[italian]{babel}
\usepackage{multirow}
\usepackage{silence}
\usepackage{tikz}
\usepackage[normalem]{ulem}
\WarningFilter{beamer}{}
\WarningFilter{metropolis}{}
\usetheme[progressbar=frametitle,titleformat=smallcaps]{metropolis}
\setbeamertemplate{frame numbering}[fraction]
\setbeamercovered{dynamic}
\definecolor{rosso}{RGB}{255, 0, 0}
\definecolor{giallo}{RGB}{254,212,23}
\hypersetup{colorlinks=true,linkcolor=black,urlcolor=rosso}
\setbeamercolor{palette primary}{fg=black, bg=giallo}
\setbeamercolor{background canvas}{bg=white}
\setbeamercolor{normal text}{fg=black}
\setbeamercolor{progress bar}{fg=rosso}
\setbeamercolor{framesubtitle}{fg=rosso}
\setbeamercolor{normal text .dimmed}{fg=giallo}
\setbeamercolor{block title alerted}{fg=rosso, bg=giallo}
\setbeamerfont{caption}{size=\tiny}
\setbeamerfont{caption name}{size=\tiny}
\setlength{\abovecaptionskip}{0pt}
\makeatletter
\metroset{block=fill}
\setlength{\metropolis@progressinheadfoot@linewidth}{1pt} 
\setlength{\metropolis@progressonsectionpage@linewidth}{1pt}
\setlength{\metropolis@titleseparator@linewidth}{1pt}
\makeatother

\title{VACANZE A BLOCCHI}
\subtitle{Gioco a coppie}
\date{}
\institute{Materiale necessario: cinque bigliettini e una penna a testa, un foglio per ciascuna coppia.}

\begin{document}

\begin{frame}[plain, noframenumbering]
    \titlepage
\end{frame}

\begin{frame}{VACANZE A BLOCCHI - CREAZIONE DEI DATI (5 MINUTI) }
    \begin{itemize}
        \item Creare cinque bigliettini, uno per ciascuno dei seguenti temi:
              \begin{enumerate}
                  \item \textbf{Luogo} visitato durante le vacanze;
                  \item \textbf{Mezzo} di trasporto utilizzato durante le vacanze;
                  \item \textbf{Persona} con la quale si è andati in vacanza;
                  \item \textbf{Attività} svolta durante le vacanze;
                  \item \textbf{Oggetto} tipico utilizzato durante le vacanze.
              \end{enumerate}
        \pause
        \item Mettere i bigliettini sulla cattedra.
    \end{itemize}
\end{frame}

\begin{frame}{VACANZE A BLOCCHI - FORMAZIONE DELLE COPPIE (5 MINUTI)}
    \begin{enumerate}
        \item Dividetevi a coppie, posizionandovi uno di fronte all'altro;
    \end{enumerate}
\end{frame}

\begin{frame}{VACANZE A BLOCCHI - ASSEGNAZIONE DEI DATI (5 MINUTI)} 
    \begin{enumerate}
        \item Ogni coppia pesca cinque bigliettini;
        \pause
        \item Ogni coppia inventa un elemento \textbf{segreto} aggiuntivo a scelta tra i cinque temi precedenti: luogo, mezzo, persona, attività, oggetto;
    \end{enumerate}
\end{frame}

\begin{frame}{VACANZE A BLOCCHI - CREAZIONE STORIA (15 MINUTI)}
    \begin{minipage}{0.98\linewidth}
        \centering
        \huge
        \textit{Ogni coppia scrive una storia sul foglio in stampatello maiuscolo 
        utilizzando tutti i dati a sua disposizione.}\\
    \end{minipage}\\
    \bigskip
    \pause
    Consegnare la storia alla cattedra, la prima coppia che consegna la storia inizia il gioco.
    \pause
    Mischiare le storie consegnate e ordinarle in un mazzo sulla cattedra.
    
\end{frame}

\begin{frame}{VACANZE A BLOCCHI - INDOVINA CHI (2 MINUTI A STORIA )}
    \begin{itemize}
        \item La coppia aggiunge il suo nome alla lavagna;
        \pause
        \item La coppia prende la prima storia dal mazzo e la legge ad alta voce;
        \pause
        \item La coppia cerca di indovinare quale coppia ha scritto la storia che sta leggendo.
    \end{itemize}
\end{frame}

\begin{frame}{VACANZE A BLOCCHI - PUNTEGGI (SEGNARLI ALLA LAVAGNA)}
    \begin{itemize}
        \item La coppia indovina di chi è la storia:
        \begin{itemize}
            \item 10 punti;
            \item Diritto a leggere la prossima storia;
            \item Può cercare di guadagnare ulteriori 5 punti indovinando l'elemento segreto della storia.
            \item La storia viene esclusa dal mazzo.
        \end{itemize}
        \item La coppia non indovina di chi è la storia: 
        \begin{itemize}
            \item Il valore della storia aumenta di 5 punti (segnarlo sul foglio);
            \item 5 punti da assegnare successivamente ai creatori della storia;
            \item Turno alla coppia indicata come autrice della storia;
            \item La storia viene messa sul fondo del mazzo.
        \end{itemize}
    \end{itemize}
\end{frame}
\end{document}